\section{Introduction}

Symbolic regression (SR) aims to recover compact, interpretable mathematical expressions from data, and has become a central problem at the intersection of scientific machine learning, scientific discovery, and automated reasoning. In recent years, the SR ecosystem has expanded rapidly: evolutionary methods, grammar-constrained search, Monte Carlo tree search variants, differentiable and neural-guided approaches, and more recently large-language-model-assisted systems have all been proposed for different parts of the SR design space. However, this algorithmic diversity has not translated into equally smooth research iteration. In practice, symbolic regression research is often slowed not only by modeling difficulty, but also by tooling fragmentation.

A typical SR method is released as a standalone repository with its own environment assumptions, data conventions, execution scripts, and result formats. Even when code is publicly available, integrating multiple methods into a common experimental workflow remains labor-intensive. The same issue appears on the data side. Different benchmark families expose different file layouts, split conventions, metadata availability, out-of-distribution (OOD) definitions, and ground-truth formula annotations. As a result, comparing methods fairly often requires substantial manual work before any scientific comparison can even begin. This creates a recurring bottleneck: a large fraction of SR research effort is spent on repository repair, environment isolation, data reshaping, evaluation alignment, and result auditing rather than on actual methodological progress.

This paper argues that symbolic regression needs a unified \emph{research substrate}: a system that makes heterogeneous methods executable under a shared protocol, makes heterogeneous datasets consumable under a shared schema, and exposes the resulting workflow to structured automation. To this end, we present \textbf{SR-Workbench}, a unified symbolic regression platform designed to reduce the cost of method onboarding, dataset canonicalization, reproducible execution, and bounded iterative improvement.

SR-Workbench is built around two normalization layers. The first is a \emph{unified method interface}. We wrap representative symbolic regression systems behind a common execution interface, isolate their dependencies through environment management, and standardize their interaction with training, prediction, and equation extraction. The current platform integrates more than ten heterogeneous SR backends spanning classical symbolic regression, search-based systems, and LLM-assisted approaches. The second is a \emph{canonical dataset schema}. We reorganize heterogeneous SR dataset families, including SRBench-, LLM-SRBench-, and SRSD-style sources, into a lightweight format with \texttt{train.csv}, \texttt{valid.csv}, \texttt{id\_test.csv}, \texttt{ood\_test.csv}, \texttt{metadata.yaml}, and optional ground-truth formulas when available. This schema is intentionally simple, but rich enough to support protocol-aware evaluation, OOD reporting, and formula-level validation.

On top of this substrate, we further introduce a modular \emph{skill layer}. Instead of treating repository integration, data conversion, protocol checking, and expression processing as ad hoc engineering chores, we package them into reusable skills. In the current version, these skills include repository onboarding for new SR tools, heterogeneous-data-to-examples conversion, expression normalization, and dataset-level validation including formula-grounded checking when ground-truth equations are available. This design makes the system easier to extend while also making the underlying workflow more explicit and auditable.

Finally, we outline and partially implement a lightweight \emph{audited evolution loop}. The loop is deliberately bounded: it starts from a runnable baseline, compiles a protocol-preserving experiment specification, executes the configured run, audits the resulting artifacts, and only then proposes limited follow-up modifications. The emphasis is not on unconstrained autonomous search, but on disciplined iteration under fixed benchmark semantics. In particular, SR-Workbench treats protocol preservation as a first-class concern: evaluation scripts, split usage, and benchmark definitions are frozen during iterative improvement, while optional ground-truth formulas are reserved for validation and analysis rather than leaked into the search process. This point is especially important in symbolic regression, where apparent gains can easily be confounded by metric drift, split leakage, inconsistent complexity definitions, or selective reporting.

Our goal in this first release is not to claim a new symbolic regression algorithm, nor to present a fully autonomous scientific agent. Instead, we target a more immediate and underdeveloped problem: the lack of a reproducible, extensible, and audit-friendly operating layer for SR research. We view SR-Workbench as a systems contribution for the SR community, one that can serve both practical experimentation and future automated research loops.

The main contributions of this work are as follows:
\begin{itemize}[leftmargin=1.5em]
    \item We build a unified symbolic regression platform that integrates more than ten heterogeneous SR tools under a common wrapper interface with isolated execution environments.
    \item We propose a lightweight canonical dataset schema that homogenizes multiple SR dataset families with explicit train/validation/ID-test/OOD-test splits, metadata, and optional ground-truth formulas.
    \item We develop a modular skill layer for symbolic regression research operations, including repository onboarding, dataset homogenization, and formula-aware dataset validation.
    \item We introduce a lightweight, auditable evolution loop for symbolic regression that supports reproducible baseline execution and bounded iterative improvement without changing benchmark semantics.
\end{itemize}

We believe this direction is timely for symbolic regression. As the field continues to expand across classical search, neural guidance, and LLM-assisted discovery, research velocity increasingly depends on whether methods and datasets can be brought into a common, trustworthy workflow. SR-Workbench is a step toward that goal.
